\subsection{1705—1722：高原危机、赋役承诺与继承}

1705年拉藏汗推翻第巴桑结嘉措的政权安排，拉萨的政教权力结构发生变动；1706年仓央嘉措被押离拉萨，其死亡或失踪的说法至今并不一致。1708年卡尔藏嘉措在青海塔尔寺附近被认定为转世之一，合法性因而同时联结青海寺院、蒙古关系和清廷册封语言。1717年准噶尔军进入拉萨、杀拉藏汗；清军1718年由四川方向入藏，在喀喇乌苏一带覆败。道路、气候、兵站与信息不足不是战报的附注，而是跨高原行动的条件。1720年清军与青海蒙古盟军进入拉萨、驱逐准噶尔力量并迎立卡尔藏嘉措为第七世达赖喇嘛；军事进入不等于一道命令已经取代西藏的政教制度。

帝国的另一端也显示出治理并非纯粹的胜负问题。1711年康熙宣布以当年丁数为定额，次年重申固定丁额、要求地方遵行。诏旨和会典能够确证国家作出的制度承诺，却不能直接换算为全国减负程度；造册、丁银、地丁并征和实际征收仍需户部题本及各省赋役档案来检验。1721年台湾朱一贵起事、一度进入府城，清廷随后派兵恢复控制，说明1684年后的跨海设治仍须面对垦殖、租佃、治安和征收的摩擦。官方“逆匪”的分类不能替代参与者的社会构成或诉求。

1722年康熙帝去世，胤禛即位。死期、即位和诏书的制度效力可由《清圣祖实录》、内阁大库遗诏档和《雍正朝起居注册》互校；围绕传位的传说不能仅因流传甚广便成为证据。新皇帝接手的是尚待处理的赋役、旗务、边疆和宗室问题。至此，现有清初编年材料完成了1644年至1722年的连续阅读；它没有授权本书跳过核验而把其后的年代写成既成事实。

\paragraph{材料位置。} 《清圣祖实录》呈现的是清廷收到信息、发出谕旨和编排凯旋的方式；藏文传统、满文奏折、寺院材料、旅行记录、地方志和赋役档案分别补足或挑战其视野。对于高原行动、丁额及台湾起事，材料所能证明的范围各不相同，不能用一类文书取代所有地方经验。

\section{定额的承诺与继承的不确定}

康熙五十年宣布以当年丁数为定额，后生人丁不再加征丁税；次年又重申固定丁额、要求地方遵行（EQ-1711-TAX-FREEZE、EQ-1712-TAX-EDICT）。《清圣祖实录》的谕旨可以确证中央规范，《大清会典事例》界定制度语言，户部题本及江苏、直隶等省的赋役档案才可能检验造册、丁银、地丁并征和实际征收。因而“滋生人丁，永不加赋”首先是一项政策承诺，不是可直接换算为全国减负程度的数字。各省执行时间、旧额、缺额和摊派的差异，要求把诏令、部议、地方账册和民众负担分开书写。它改变了丁额增长的制度预期，并为雍正时期的赋役调整留下条件，却没有消灭地方财政压力。

康熙六十一年十一月，康熙帝去世，皇四子胤禛即位（EQ-1722-KANGXI-DEATH）。死期、即位与诏书的制度效力，可由《清圣祖实录》、内阁大库遗诏档和《雍正朝起居注册》互校。《清史稿》提供了后出的编年索引，却不能使传位过程中的每一个传说自动成立。康熙晚年缺少公开稳定的继承机制，宗室、旗人和官僚围绕信息形成竞争和猜测；这一背景可以说明何以出现争议，不能把争议本身当成阴谋的证据。新皇帝继承的，是尚待处理的财政、旗务、边疆和宗室问题，而非一套已经没有摩擦的盛世机器。


\subsection{事件链与材料限度}
赋役、继承与地方执行之间的联系，须在具体账册、诏令和社会实践中逐项追踪。

\paragraph{康熙五十年（1711年）：康熙帝宣布以康熙五十年丁数为定额，后生人丁不再加征丁税}
\eventanchor{EQ-1711-TAX-FREEZE}{康熙帝宣布以康熙五十年丁数为定额，后生人丁不再加征丁税}。本节点叙述该事件本身。清在京师与多地军政线路相接的尺度的处境首先受长期战争后清廷寻求稳定赋役预期，人口登记与丁税征收的摩擦已显现制约。《清圣祖实录》康熙五十年丁银谕；《清史稿·食货志》；户部档案与《大清会典事例》赋役条留下可核的线索；可辨的后果是政策改变丁额增长的制度预期，为雍正朝摊丁入亩等改革留下条件；实际减负须按省县账册验证。原有置信注记：诏令可靠；各省执行时间、丁银与地丁并征口径和实际负担不能由诏令推定

\paragraph{康熙五十年（1711年）：康熙帝宣布以康熙五十年丁数为定额，后生人丁不再加征丁税}
\eventanchor{EQ-1711-TAX-FREEZE-AFTERMATH}{康熙帝宣布以康熙五十年丁数为定额，后生人丁不再加征丁税} 置于京师与多地军政线路相接的尺度中阅读，本节点追踪「康熙帝宣布以康熙五十年丁数为定额，后生人丁不再加征丁税」之后可见的余波。长期战争后清廷寻求稳定赋役预期，人口登记与丁税征收的摩擦已显现构成行动者清必须面对的条件。取证从《清圣祖实录》康熙五十年丁银谕；《清史稿·食货志》；户部档案与《大清会典事例》赋役条出发，因而只能把变化谨慎地连到政策改变丁额增长的制度预期，为雍正朝摊丁入亩等改革留下条件；实际减负须按省县账册验证。原有置信注记：诏令可靠；各省执行时间、丁银与地丁并征口径和实际负担不能由诏令推定

\paragraph{康熙五十年（1711年）：康熙帝宣布以康熙五十年丁数为定额，后生人丁不再加征丁税}
\eventanchor{EQ-1711-TAX-FREEZE-DECISION}{康熙帝宣布以康熙五十年丁数为定额，后生人丁不再加征丁税} 标记本节点定位围绕「康熙帝宣布以康熙五十年丁数为定额，后生人丁不再加征丁税」形成的处置与调度记录，涉及清。京师与多地军政线路相接的尺度不是抽象背景：长期战争后清廷寻求稳定赋役预期，人口登记与丁税征收的摩擦已显现。《清圣祖实录》康熙五十年丁银谕；《清史稿·食货志》；户部档案与《大清会典事例》赋役条所见的顺序随后指向政策改变丁额增长的制度预期，为雍正朝摊丁入亩等改革留下条件；实际减负须按省县账册验证，但原有置信注记：诏令可靠；各省执行时间、丁银与地丁并征口径和实际负担不能由诏令推定

\paragraph{康熙五十年（1711年）：康熙帝宣布以康熙五十年丁数为定额，后生人丁不再加征丁税}
在京师与多地军政线路相接的尺度，\eventanchor{EQ-1711-TAX-FREEZE-PRELUDE}{康熙帝宣布以康熙五十年丁数为定额，后生人丁不再加征丁税} 让清的一个选择或压力有了可查的坐标。本节点回看「康熙帝宣布以康熙五十年丁数为定额，后生人丁不再加征丁税」之前已可见的条件。前因可写为长期战争后清廷寻求稳定赋役预期，人口登记与丁税征收的摩擦已显现；《清圣祖实录》康熙五十年丁银谕；《清史稿·食货志》；户部档案与《大清会典事例》赋役条限定了记录，后续变化是政策改变丁额增长的制度预期，为雍正朝摊丁入亩等改革留下条件；实际减负须按省县账册验证。原有置信注记：诏令可靠；各省执行时间、丁银与地丁并征口径和实际负担不能由诏令推定

\paragraph{康熙五十一年（1712年）：清廷再次谕令固定丁额并要求地方遵行，围绕丁银征收与造册形成持续行政问题}
\eventanchor{EQ-1712-TAX-EDICT}{清廷再次谕令固定丁额并要求地方遵行，围绕丁银征收与造册形成持续行政问题} 所关联的是清在省、府县与地方社会相互牵动的尺度中的行动。本节点叙述该事件本身。它从上一年诏令需经户部、督抚和州县造册落实，旧有缺额、摊派和财政压力造成执行摩擦展开，借《清圣祖实录》康熙五十一年赋役谕；《清史稿·食货志》；户部题本与江苏、直隶等省赋役档进入可检验的年序，并以“永不加赋”成为政策语言与后续改革依据，但登记人口、实际税负和地区差异不能被单一口号遮蔽影响下一段安排。原有置信注记：重复申谕可靠；它证明中央规范而非各地立即停止加派，需追踪地方账册与奏报

\paragraph{康熙五十一年（1712年）：清廷再次谕令固定丁额并要求地方遵行，围绕丁银征收与造册形成持续行政问题}
省、府县与地方社会相互牵动的尺度中的\eventanchor{EQ-1712-TAX-EDICT-AFTERMATH}{清廷再次谕令固定丁额并要求地方遵行，围绕丁银征收与造册形成持续行政问题} 不能离开清来解释。本节点追踪「清廷再次谕令固定丁额并要求地方遵行，围绕丁银征收与造册形成持续行政问题」之后可见的余波。上一年诏令需经户部、督抚和州县造册落实，旧有缺额、摊派和财政压力造成执行摩擦说明其形成的条件；《清圣祖实录》康熙五十一年赋役谕；《清史稿·食货志》；户部题本与江苏、直隶等省赋役档提供来源路径，“永不加赋”成为政策语言与后续改革依据，但登记人口、实际税负和地区差异不能被单一口号遮蔽则是材料能够支持的近时结果。原有置信注记：重复申谕可靠；它证明中央规范而非各地立即停止加派，需追踪地方账册与奏报

\paragraph{康熙五十一年（1712年）：清廷再次谕令固定丁额并要求地方遵行，围绕丁银征收与造册形成持续行政问题}
\eventanchor{EQ-1712-TAX-EDICT-DECISION}{清廷再次谕令固定丁额并要求地方遵行，围绕丁银征收与造册形成持续行政问题} 把清放回省、府县与地方社会相互牵动的尺度的道路、城镇、边界或制度网络。本节点定位围绕「清廷再次谕令固定丁额并要求地方遵行，围绕丁银征收与造册形成持续行政问题」形成的处置与调度记录。上一年诏令需经户部、督抚和州县造册落实，旧有缺额、摊派和财政压力造成执行摩擦。本段采用《清圣祖实录》康熙五十一年赋役谕；《清史稿·食货志》；户部题本与江苏、直隶等省赋役档核对其顺序，随后可追到“永不加赋”成为政策语言与后续改革依据，但登记人口、实际税负和地区差异不能被单一口号遮蔽。原有置信注记：重复申谕可靠；它证明中央规范而非各地立即停止加派，需追踪地方账册与奏报

\paragraph{康熙五十一年（1712年）：清廷再次谕令固定丁额并要求地方遵行，围绕丁银征收与造册形成持续行政问题}
本节点回看「清廷再次谕令固定丁额并要求地方遵行，围绕丁银征收与造册形成持续行政问题」之前已可见的条件，故\eventanchor{EQ-1712-TAX-EDICT-PRELUDE}{清廷再次谕令固定丁额并要求地方遵行，围绕丁银征收与造册形成持续行政问题} 不只是一个年份标签。清在省、府县与地方社会相互牵动的尺度承受的压力来自上一年诏令需经户部、督抚和州县造册落实，旧有缺额、摊派和财政压力造成执行摩擦。《清圣祖实录》康熙五十一年赋役谕；《清史稿·食货志》；户部题本与江苏、直隶等省赋役档可回查该节点；“永不加赋”成为政策语言与后续改革依据，但登记人口、实际税负和地区差异不能被单一口号遮蔽是其进入后续年序的方式。原有置信注记：重复申谕可靠；它证明中央规范而非各地立即停止加派，需追踪地方账册与奏报

\paragraph{康熙六十一年十一月（1722年）：康熙帝去世，皇四子胤禛即位为雍正帝，早期清朝的长期统治转入新的继承阶段}
\eventanchor{EQ-1722-KANGXI-DEATH}{康熙帝去世，皇四子胤禛即位为雍正帝，早期清朝的长期统治转入新的继承阶段} 所见的行动者为清，空间尺度为京师与多地军政线路相接的尺度。本节点叙述该事件本身。康熙晚年未公开稳定继承机制，宗室、旗人和官僚围绕皇位信息形成竞争与猜测使它成为具体事件链的一环；《清圣祖实录》康熙六十一年十一月；《清史稿·圣祖本纪、世宗本纪》；《雍正朝起居注册》与内阁大库遗诏档与雍正即位后将处理财政、旗务、边疆和宗室问题；传位细节不能以阴谋叙事代替可核文书。原有置信注记：去世与即位日期可靠；遗诏、传位过程和储位争议受政治清算与后世传说影响，须以档案版本互校

\paragraph{康熙六十一年十一月（1722年）：康熙帝去世，皇四子胤禛即位为雍正帝，早期清朝的长期统治转入新的继承阶段}
从京师与多地军政线路相接的尺度看，\eventanchor{EQ-1722-KANGXI-DEATH-AFTERMATH}{康熙帝去世，皇四子胤禛即位为雍正帝，早期清朝的长期统治转入新的继承阶段} 留下本节点追踪「康熙帝去世，皇四子胤禛即位为雍正帝，早期清朝的长期统治转入新的继承阶段」之后可见的余波的材料坐标。清所面对的条件是康熙晚年未公开稳定继承机制，宗室、旗人和官僚围绕皇位信息形成竞争与猜测。《清圣祖实录》康熙六十一年十一月；《清史稿·圣祖本纪、世宗本纪》；《雍正朝起居注册》与内阁大库遗诏档支持的结果为雍正即位后将处理财政、旗务、边疆和宗室问题；传位细节不能以阴谋叙事代替可核文书。原有置信注记：去世与即位日期可靠；遗诏、传位过程和储位争议受政治清算与后世传说影响，须以档案版本互校

\paragraph{康熙六十一年十一月（1722年）：康熙帝去世，皇四子胤禛即位为雍正帝，早期清朝的长期统治转入新的继承阶段}
\eventanchor{EQ-1722-KANGXI-DEATH-DECISION}{康熙帝去世，皇四子胤禛即位为雍正帝，早期清朝的长期统治转入新的继承阶段} 记录清在京师与多地军政线路相接的尺度的一次变化。本节点定位围绕「康熙帝去世，皇四子胤禛即位为雍正帝，早期清朝的长期统治转入新的继承阶段」形成的处置与调度记录。其发生与康熙晚年未公开稳定继承机制，宗室、旗人和官僚围绕皇位信息形成竞争与猜测相连；《清圣祖实录》康熙六十一年十一月；《清史稿·圣祖本纪、世宗本纪》；《雍正朝起居注册》与内阁大库遗诏档把事件系回来源，雍正即位后将处理财政、旗务、边疆和宗室问题；传位细节不能以阴谋叙事代替可核文书显示压力如何向后转移。原有置信注记：去世与即位日期可靠；遗诏、传位过程和储位争议受政治清算与后世传说影响，须以档案版本互校

\paragraph{康熙六十一年十一月（1722年）：康熙帝去世，皇四子胤禛即位为雍正帝，早期清朝的长期统治转入新的继承阶段}
\eventanchor{EQ-1722-KANGXI-DEATH-PRELUDE}{康熙帝去世，皇四子胤禛即位为雍正帝，早期清朝的长期统治转入新的继承阶段} 的地点与行动范围属于京师与多地军政线路相接的尺度，行动者是清。本节点回看「康熙帝去世，皇四子胤禛即位为雍正帝，早期清朝的长期统治转入新的继承阶段」之前已可见的条件。康熙晚年未公开稳定继承机制，宗室、旗人和官僚围绕皇位信息形成竞争与猜测。读《清圣祖实录》康熙六十一年十一月；《清史稿·圣祖本纪、世宗本纪》；《雍正朝起居注册》与内阁大库遗诏档时可见其记录边界，最小后果只能写作雍正即位后将处理财政、旗务、边疆和宗室问题；传位细节不能以阴谋叙事代替可核文书。原有置信注记：去世与即位日期可靠；遗诏、传位过程和储位争议受政治清算与后世传说影响，须以档案版本互校
